% ===========================================================================
% sections_v2/contributions.tex -- round-8 rewrite: a qualitative list only.
% Every number in this paper lives in Sec. VII and in the % source: comments
% there.  This file states what the paper contributes and what it does not
% claim, with section pointers and no quantitative claim of its own, so that
% the introduction can never drift from the results again.
% Cross-file labels: sec:rw-concurrent is defined in sections_v2/related_work.tex.
% ===========================================================================
\subsection{Contributions}
\label{sec:contributions}

\begin{enumerate}

\item \textbf{A physically closed relaxation for discrete structural decisions
  inside a production water-network solver, with an adjoint that matches central
  differences to the finite-difference floor ($3\times10^{-10}$ at interior points,
  $1.4\times10^{-3}$ at one-hot closed-branch coordinates)}
  (\cref{sec:method,sec:e0lc}).  The relaxation lives inside the solver's own
  friction law, so every probability vector over design tokens is a network the
  unmodified solver can simulate.  Closure is a measured neighbourhood property:
  the floor below which unchosen branches vanish is swept, and its price
  reported at every value.

\item \textbf{The bias of a first-order pathwise gradient depends on the
  \emph{type} of decision it is asked about} (\cref{sec:e0lc,sec:e1lc}).  On
  the production solver the same adjoint ranks catalogue edits faithfully and
  structural re-assignments poorly, and no estimator in the compared family
  repairs the asymmetry.

\item \textbf{A measurement protocol for post-training design agents, and the
  accounting it produces} (\cref{sec:e2,sec:e2B} and Supplementary Sec.~S-VII).  Paired per-task
  tests at the stated seed count, an endpoint-versus-best-checkpoint census over
  every archived run, and a decomposition of each objective reduction into
  design and penalty parts.  Applied here it shows post-training buying
  constraint satisfaction by spending design quality on one benchmark, and
  regressing against its own warm start on the other.  The protocol's one
  unqualified positive is at matched budget on a 6{,}599-parameter factorised
  policy, where pathwise halves scalar-reward regret (\cref{sec:baselines}).
  % source: ddl/case/analysis/v2A_table/loggers/20260915T175120Z_gpuserver/
  %         results_v2A_table.json -> table, paired_wilcoxon; break-even from
  %         ddl/case/analysis/consistency_inputs/loggers/20260915T154717Z_gpuserver/
  %         results_consistency.json -> c_breakeven.
  The break-even arithmetic of \cref{sec:discussion} states the amortisation
  threshold as a falsifiable prediction for the minutes-per-solve regime.

\item \textbf{A protocol for scoring gradient estimators when the artefact
  that ships is a hard design} (\cref{sec:e1lc,sec:discussion}).  The
  conventional metric, cosine against a high-sample unbiased reference, is only
  a description of a nine-arm table here (\cref{sec:e1lc}), so we do not use it.
  The substitute is discrete edit selection, validated by a deterministic probe:
  paired over ten confident policy points the exact adjoint beats the
  scalar-reward arm (\cref{sec:e1lc}, one-sided $p=0.0020$).
  % source: ddl/case/analysis/round10_reanalysis/loggers/20260915T201839Z_gpuserver/
  %         results_round10.json (rho, permutation p); edit-selection paired test from
  %         ddl/case/benchB/E1d_edit_ranking/loggers/20260915T150207Z_edit5seed.

\item \textbf{Benchmark validity as a result} (\cref{sec:setup} and
  Supplementary Sec.~S-X).  The first benchmark of this line of work failed a
  physicality audit that its own labels made computable.  The repair is verified
  against enumeration wherever enumeration is possible, and every result carries
  the physics configuration it was run under.

\end{enumerate}

\subsection{What we do not claim}
\label{sec:notclaim}

Four negative statements bound the five positive ones.  We do not claim that
pathwise estimators beat score-function estimators in general; our evidence is
about decision types, not estimators.  We do not claim that post-training
improves design quality: on both benchmarks it either buys feasibility at a
design cost or regresses against its warm start.  We do not claim a scale
result, since no 27B run is reported here
(\cref{sec:e27b}), nor that a language model beats constrained search where
search is cheap (\cref{sec:baselines}).  A concurrent paper made the Global
Gradient Method differentiable before this draft froze \cite{mu2026gga}; our
relaxation and the weight-level use of its adjoint are independent of that work
(\cref{sec:rw-concurrent}).
